\chapter{Requisitos do Projeto}
\label{chap:requisitos_final}

\orientacao{Tabela consolidada com a trajetória completa do projeto, uma coluna por macroentrega, do M01 ao final. Ela serve de registro de execução para o parceiro e para a auditoria. Preencha a coluna \textbf{Verificação} nas linhas \parc{} e \done{}, com o método e o ponteiro para a evidência no relatório. Para o uso detalhado da tabela, veja o capítulo de requisitos de \texttt{main.tex}.}

Este capítulo consolida os requisitos do projeto \ph{NOME DO PROJETO} e o status de cada um ao final de cada macroentrega. O status é indicado pelas siglas \todo{} (não iniciado), \wip{} (em andamento), \parc{} (parcialmente implementado) e \done{} (concluído e verificado).

\section{Requisitos Funcionais}

\begin{tabelarequisitos}{M01,M02,M03,M04}
\req{RF-XXX-01}{\ph{O sistema deve ...}}{\wip & \wip & \done & \done}{Ensaio, Tab.~\ref{tab:wp1_resultados}}
\req{RF-XXX-02}{\ph{O sistema deve ...}}{\todo & \wip & \parc & \done}{Demonstração, Fig.~\ref{fig:wp1}}
\end{tabelarequisitos}

\section{Requisitos Não Funcionais}

\begin{tabelarequisitos}{M01,M02,M03,M04}
\req{RNF-XXX-01}{\ph{O sistema deve ...}}{\todo & \wip & \wip & \done}{Análise, Anexo~\ref{anx:a}}
\end{tabelarequisitos}

\section{Leitura da Tabela}

Dos \ph{N} requisitos do projeto, \ph{N} estão concluídos e verificados ao final da última macroentrega. \ph{Onde o avanço se concentrou, o que permanece em aberto e por qual dependência.}
